\documentclass[11pt]{article}
\usepackage{amsmath, amssymb}
\usepackage{geometry}
\usepackage{array}
\usepackage{booktabs}
\geometry{margin=0.7in}
\setlength{\parindent}{0pt}
\setlength{\parskip}{0.35em}
\renewcommand{\arraystretch}{1.15}

\newcommand{\cheatsheettitle}[2]{%
\begin{center}
{\LARGE #1}\\[0.4em]
{\large #2}
\end{center}
}


\begin{document}
\cheatsheettitle{Algebra II Radicals and Rational Exponents Cheatsheet}{Module 6}

\section*{Rational Exponent Forms}
\[
a^{1/n}=\sqrt[n]{a}
\]
\[
a^{m/n}=\sqrt[n]{a^m}=(\sqrt[n]{a})^m
\]
\[
a^{-m/n}=\frac{1}{a^{m/n}}
\]

\section*{Radical Rules}
\[
\sqrt[n]{ab}=\sqrt[n]{a}\sqrt[n]{b}
\]
\[
\sqrt[n]{\frac ab}=\frac{\sqrt[n]{a}}{\sqrt[n]{b}}\qquad (b\ne0)
\]
\[
\sqrt[n]{a^n}=
\begin{cases}
|a|, & n\text{ even}\\
a, & n\text{ odd}
\end{cases}
\]

\section*{Simplifying}
Pull out perfect powers.
\[
\sqrt[3]{54x^5}=\sqrt[3]{27x^3\cdot2x^2}=3x\sqrt[3]{2x^2}
\]
Combine only like radicals:
\[
4\sqrt7-9\sqrt7=-5\sqrt7
\]

\section*{Rationalizing}
Single radical denominator:
\[
\frac{5}{\sqrt3}=\frac{5\sqrt3}{3}
\]
Binomial radical denominator:
\[
\frac{1}{a+\sqrt b}\cdot\frac{a-\sqrt b}{a-\sqrt b}
=\frac{a-\sqrt b}{a^2-b}
\]

\section*{Solving Radical Equations}
\begin{enumerate}
  \item Isolate one radical.
  \item Raise both sides to the needed power.
  \item Solve the resulting equation.
  \item Check in the original equation.
\end{enumerate}
Example template:
\[
\sqrt{A}=B\Rightarrow A=B^2
\]
with
\[
B\ge0
\]

\section*{Extraneous Solutions}
Squaring can introduce false solutions. Every solution must satisfy the original equation and any domain restrictions.

\section*{Quick Checks}
\begin{itemize}
  \item Even-index radicals require nonnegative radicands over the reals.
  \item Odd-index radicals can have negative radicands.
  \item Rational exponents follow exponent rules.
  \item Simplified radicals have no removable perfect-power factors.
\end{itemize}

\end{document}
